% [原文 p.267]
\par\medskip\noindent\textbf{7.21}（证明的步骤与记号）\quad 设 $f'_{tf}\colon X'_{tf}\to C'$ 是诱导态射，$B'_n,B'_{tf}$ 分别是 $B$ 向 $X'_n,X'_{tf}$ 的拉回。由命题 5.15 与 7.21(1)，$X'_n$ 是 Q-分解的。特别地，下面所构造的 $\operatorname{Ex}(X'_{tf}\to X'_n)$ 是纯余维 1 的（推论 2.63）。另外，往 $B$ 中加入纤维没有任何影响，故可假设 $B$ 包含所有奇异纤维。于是 $f$ 在 $X\setminus B$ 上光滑，且 $X'_n\setminus B'_n$ 非奇异。由此 $X'_{tf}\setminus B'_{tf}\simeq X'_n\setminus B'_n$。由于 $\operatorname{Ex}(X'_{tf}\to X'_n)\subset B'_{tf}$，为证定理 7.17(1)，只需确认对每个 $c'\in C'$，$(f'_{tf})^*(c')\cup B'_{tf}$ 是 snc 的。

定理 7.19 按以下步骤证明。
\begin{enumerate}[label=(\arabic*),leftmargin=2.4em,itemsep=1pt,topsep=2pt,parsep=0pt]
\item 确认 $X'_n$ 只有标准商奇点（命题 7.23）。
\item 设 $X'_t\to X'_n$ 是标准终端化（注意 6.24）。
\item 记 $q_t\colon X'_t\to X$，并把 $\operatorname{red}(q_t^{-1}(B))=\sum_t B'^{(\ell)}_t$ 表示出来。其中 $B'^{(\ell)}_t$ 是由以下条件定义的约化除子：
\[
G_1,G_2\ \text{属于同一个}\ B'^{(\ell)}_t
\iff
\text{对各除子}\ E\subset B,\ \operatorname{mult}_{G_1}(q_t^*E)=\operatorname{mult}_{G_2}(q_t^*E).
\]
由定理 6.32，按任意顺序逐个把 $B'^{(\ell)}_t$ 变成 Q-Cartier 除子。这样便得到终端化 $X'_{tf}\to X'_t$，使双有理变换 $B'^{(\ell)}_{tf}$ 全都是 Q-Cartier 的。
\item 确认 $f'_{tf}\colon X'_{tf}\to C'$ 的所有纤维都是约化的（引理 7.24）。
\item 确认 $X'_{tf}$ 与 $B'^{(\ell)}_{tf}$ 全都非奇异。
\item 确认 $B'_{tf}$ 在局部解析意义下是 nc 除子。\hfill$\square$
\end{enumerate}

% [Tutor Note: 阶段 3 把重数相同的素除子归为一组，是为了让之后逐个 Q-Cartier 化的爆破保持等变，从而把一般情形归约到命题 7.23 的单项式模型。]
上述阶段 3 中引入 $B'^{(\ell)}_t$ 的理由由下述引理说明。

\begin{lemmap}{7.22}
要执行 7.21 的各阶段，只需在命题 7.23 中令 $m=3$ 所得的代数芽 $f(x_1,x_2,x_3)=\prod x_i^{a_i}$、$p(t)=t^d$ 的情形进行计算即可。
\end{lemmap}

\begin{proof}
在 $X\setminus f^{-1}(c)$ 上无需证明。用 $W\subset X$ 表示点 $\xi\in f^{-1}(c)$ 的充分小的解析邻域。此时 $X'|_W=W'$、$X'_n|_W=W'_n$ 成立。又由注意 6.24，得到 $X'_t|_W=W'_t$。由于 $B$ 是 snc 除子，$B$ 的任意不可约分支在 $W$ 上的限制都是 $B|_W$ 的不可约分支。由阶段 3

% [原文 p.268]
的构造，任意 $B'^{(\ell)}_t$ 的限制都与某个 $(B|_W)'^{(\ell)}_t$ 一致，或者是空集。由此，若把阶段 3 中爆破的顺序选得与之对应，则 $X'_{tf}|_W=W'_{tf}$（注意 6.33）。因为目的只是确认 $X'_{tf}$ 的局部解析性质，所以确认 $W'_{tf}$ 的相应性质也无妨。即使把 $\xi\in X$、$\xi\in B$、$\xi$ 处的 $f$、$p^{-1}(c)$ 处的 $p$ 等芽的组替换为解析同型的其他芽的组，$W'_{tf}$ 仍解析同型（注意 6.24、注意 6.33）。
\end{proof}

阶段 1 由下述命题处理。

\begin{propositionp}{7.23}
设 $X=(0\in\C^m)$ 是芽，$B$ 是满足 $\{\prod x_i^{a_i}=0\}\subset B\subset\{\prod x_i=0\}$ 的除子。其中 $a_i\ge0$ 是满足 $\sum a_i>0$ 的整数，$d>0$ 是所有非零 $a_i$ 的公倍数，$q_n\colon X'_n\to X$ 是射影。其中 $X'_n$ 是
\[
X':=\Bigl(\prod x_i^{a_i}=t^d\Bigr)\subset(0\in X\times\C^1)
\]
的正规化，$f'_n\colon X'_n\to X'\to\C^1$ 是第 2 投影，记 $B'_n=q_n^*B$。此时：
\begin{enumerate}[label=(\arabic*),leftmargin=2.4em,itemsep=1pt,topsep=2pt,parsep=0pt]
\item $K_{X'_n}$ 是 Cartier 的，且 $D'_n:=((f'_n)^*t=0)$ 是约化的。
\item $X'_n\setminus D'_n$ 非奇异，且 $X'_n$ 只有商奇点。
\item $K_{X'_n}+D'_n=q_n^*(K_X+\operatorname{red}(q_n(D'_n)))$。
\item $(X'_n,0)$ 是标准的，且 $(X'_n,D'_n)$ 是 lc 的。
\end{enumerate}
\end{propositionp}

\begin{proof}
令 $a_i=0$ 时 $b_i=1$，$a_i\ne0$ 时 $b_i=d/a_i$。对 1 的任意 $d$ 次根 $\eps$，用下式定义 $\pi\colon(0\in\C^m)\to(0\in X')$：
\[
\pi(u_1,\cdots,u_m)=\Bigl(u_1^{b_1},\cdots,u_m^{b_m},\eps\prod_{a_i>0}u_i\Bigr).
\]
此时，$\pi$ 定义到 $X'_n$ 的不可约分支的芽上的有限满射 $\pi_n$。又若让 $\eps$ 变动，则 $\operatorname{im}\pi_n$ 覆盖 $X'_n$。由于
\[
\pi_n^*(f'_n)^*t=\pi^*t=\eps\prod_{a_i>0}u_i,
\]
$D'_n$ 是约化的，且 $\pi_n$ 在余维 1 处是 étale 的。由适当的有限群 $G$ 有 $\operatorname{im}\pi_n\simeq\C^m/G$，故 $X'_n$ 只有商奇点。

由于 $q_n\colon X'_n\to X$ 在 $X'_n\setminus D'_n$ 上非分歧，可知 $X'_n\setminus D'_n$ 非奇异。这就得到 (2)。又由于 $\pi_n$ 在余维 1 处 étale，(3) 由显然的
\[
K_{\C^m}+(\pi^*t=0)=(pr_1\circ\pi)^*(K_X+\operatorname{red}(q_n(D'_n)))
\]
得到。从而

% [原文 p.269]
$K_{X'_n}$ 是 Cartier 的。又由推论 2.31，$(X,\operatorname{red}(q_n(D'_n)))$ 是 lc 的，故由命题 5.20，$(X'_n,D'_n)$ 是 lc 的。为证明 $(X'_n,0)$ 是标准的，设 $E$ 是 $X'_n$ 上方的任意例外除子。若 $\operatorname{center}_{X'_n}(E)\not\subset D'_n$，则由 (2) 有 $a(E,X'_n)>0$。若 $\operatorname{center}_{X'_n}(E)\subset D'_n$，则由 $D'_n$ 是 Cartier 的且 $(X'_n,D'_n)$ 是 lc 的，得到
\[
a(E,X'_n)\ge a(E,X'_n,D'_n)+1\ge0.
\]
\end{proof}

7.21 的阶段 2、3 是自动的。下面的引理完成阶段 4 的确认。

\begin{lemmap}{7.24}
设 $Y$ 是标准的，$D\subset Y$ 是 Cartier 除子，且假设 $(Y,D)$ 是 lc 的。设 $g\colon Y'\to Y$ 是任意 crepant 双有理态射。此时，$g^*D\subset Y'$ 是约化的。
\end{lemmap}

\begin{proof}
记 $g^*D=\sum b_iD_i$。由于 $K_{Y'}\equiv g^*K_Y$，得到
\[
K_{Y'}\equiv g^*(K_Y+D)+\sum(-b_i)D_i.
\]
由于 $(Y,D)$ 是 lc 的，对每个 $i$ 有 $b_i\le1$，故 $g^*D$ 是约化的。
\end{proof}

（群作用）下面证明 $X'_{tf}$ 非奇异。这不是由一般理论形式地推出的事实，必须使用 $(\C^*)^3$-作用。不过，为避免系统性的准备，这里只用其极大挠子群 $\mu_\infty$（由 1 的所有单位根构成的群）。（处理 $(\C^*)^3$ 时，必须把它作为拓扑群或代数群来处理。）这里先指出：$\mu_\infty$ 是可除（divisible）群，且 $\mu_r$（由 1 的所有 $r$ 次根构成的群）是 $\mu_\infty$ 的全部有限子群。

% [Tutor Note: 这里绕开 $(\C^*)^3$ 的代数群或拓扑群理论，只用其挠子群 $\mu_\infty$：可除群在有限集上的作用平凡，足以保证不动点与等变性。]
\begin{lemma}{7.25}
按如下方式定义 $(\mu_\infty)^m$ 对命题 7.23 中引入的 $X'$ 的作用：
\[
(\lambda_1,\cdots,\lambda_m)\colon(x_1,\cdots,x_m,t)\mapsto(\lambda_1^dx_1,\cdots,\lambda_m^dx_m,t\prod\lambda_i^{a_i}).
\]
这一作用可延拓为 $X'_n$ 上的、具有有限核的 $(\mu_\infty)^m$-作用，使 $B'_n$ 是 $(\mu_\infty)^m$-不变的（即同态 $(\mu_\infty)^m\to\operatorname{Aut}(X'_n)$ 的核是有限群）。
\end{lemma}

\begin{corollaryp}{7.26}
由引理 7.25 定义的 $X'_n$ 上的 $(\mu_\infty)^3$-作用，可延拓为 $X'_{tf}$ 上的、具有有限核的 $(\mu_\infty)^3$-作用，使 $B'^{(\ell)}_{tf}$ 是 $(\mu_\infty)^3$-不变的。
\end{corollaryp}

\begin{proof}
$(\mu_\infty)^3$ 作用于标准终端化 $X'_t$（注意 6.24），且 $B'_t$ 是 $(\mu_\infty)^3$-不变的。由于 $(\mu_\infty)^3$ 可除，其对有限集的作用都是平凡的。因此，

% [原文 p.270]
$B'_t$ 的每个不可约分支都不变，且 $X'_{tf}\to X'_t$ 途中的爆破是等变的（equivariant）。
\end{proof}

$X'_{tf}$ 的奇点是终端的，且带有 $(\mu_\infty)^3$-作用。不过存在带有 $(\C^*)^3$-作用的指数 1 终端奇点（例如 $(x_1x_2=x_3x_4)\subset\C^4$），这些被关于 $B'^{(\ell)}_{tf}$ 的条件（推论 7.30）排除。下面叙述的是关键的引理。

\begin{lemmap}{7.27}
对任意不可约分支 $E\subset B'_{tf}$，$T_E$ 是由下式给出的 $(\mu_\infty)^3$ 的子群：
\[
T_E=\{(t^{v_E(x_1)},t^{v_E(x_2)},t^{v_E(x_3)})\mid t\in\mu_\infty\}\subset(\mu_\infty)^3.
\]
其中 $v_E(\ )$ 表示沿 $E$ 的赋值。此时，$T_E\simeq\mu_\infty$，且 $E$ 的每个点都被 $T_E$ 固定。
\end{lemmap}

\begin{proof}
由下式定义 $(\mu_\infty)^3$-等变的有理映射 $h\colon E\dashrightarrow\PP^2$：
\[
X\ni(x_1,x_2,x_3)\mapsto(x_1^{v_2v_3}:x_2^{v_3v_1}:x_3^{v_1v_2}),
\]
其中 $v_i=v_E(x_i)$，而 $\PP^2$ 上的 $(\mu_\infty)^3$-作用由
\[
(\lambda_1,\lambda_2,\lambda_3)\colon(y_1:y_2:y_3)\mapsto(\lambda_1^{d\cdot v_2v_3}y_1:\lambda_2^{d\cdot v_3v_1}y_2:\lambda_3^{d\cdot v_1v_2}y_3)
\]
给出。

$\PP^2$ 的 $(\mu_\infty)^3$-不变闭不可约 Zariski 子集只有 3 点 $(1:0:0),(0:1:0),(0:0:1)$、3 条直线 $(y_i=0)$ 以及 $\PP^2$ 本身。

指数选得使 3 个坐标函数沿 $E$ 的赋值相等。因此，$h(E)$ 不含于 $\PP^2$ 的坐标直线中。$h(E)$ 的闭包是 $\PP^2$ 的 $(\mu_\infty)^3$-不变子集，故 $h$ 是支配的。

注意 $T_E$ 在 $\PP^2$ 上平凡地作用。因此 $T_E$ 作用于 $h$ 的一般纤维，而那是有限集。由于 $T_E$ 可除，$T_E$ 在 $E$ 上的作用是平凡的。
\end{proof}

下面在更一般的设定下确认 $B'^{(\ell)}_{tf}$、$X'_{tf}$ 的非奇异性。

\begin{lemmap}{7.28}
设 $(0\in Y)$ 是三维超曲面的芽，且 $Y$ 上有具有有限核的 $\mu_\infty$-作用。再假设 $Y\setminus\{0\}$ 非奇异，且存在使每点都被 $\mu_\infty$ 固定的 Weil 除子 $0\in E\subset Y$。此时，$E$ 非奇异。此外，$E$ 是 Q-Cartier 当且仅当 $Y$ 非奇异。
\end{lemmap}

% [原文 p.271]
\begin{proof}
存在使 $\mu_a=\ker[\mu_\infty\to\operatorname{Aut}Y]$ 的 $a>0$。因此，由 $\mu_\infty\simeq\mu_\infty/\mu_a$，可以假设作用是忠实的，即 $\mu_\infty\subset\operatorname{Aut}Y$。

取 $r\gg1$，设 $\mu_r\subset\mu_\infty$。由引理 7.29，把 $Y\subset\C^4$ 嵌入，使得
\[
\lambda\colon(y_1,\cdots,y_4,\phi)\mapsto(\lambda^{c_1}y_1,\cdots,\lambda^{c_4}y_4,\lambda^e\phi)\qquad(\lambda\in\mu_r,\ c_i,e\in\Z)
\]
是 $\mu_r$-作用。其中 $\phi$ 是 $Y$ 的方程。可以对适当的 $k$ 假设 $c_i\equiv0\ (r)\iff i\le k$。$\mu_r$-作用的不动点集是 $(y_{k+1}=\cdots=y_4=0)$。由于其维数至少为 2，故 $k\in\{4,3,2\}$。首先，由于 $\mu_\infty$-作用是忠实的，$k=4$ 不可能。

若 $k=3$，则 $\phi\in(y_4)$ 或 $\phi\in\C\{y_1,y_2,y_3,y_4^r\}$。但作用忠实，故只能出现后者。这里令 $\phi_i=\partial_{y_i}\phi$，则由理想 $(\phi,\phi_1,\cdots,\phi_4)$ 定义的特异点概形 $\operatorname{Sing}_{\mathrm{sch}}(Y)\subset Y$ 与坐标系 $y_1,\cdots,y_4$、方程 $\phi$ 以及 $r$ 的选取无关。若 $Y$ 是奇异的，则 $\operatorname{Sing}_{\mathrm{sch}}(Y)$ 是 Artin 概形 $\ni0$，且容易看出 $(\phi,\phi_1,\cdots,\phi_4)\subset(y_1,\cdots,y_3,y_4^{r-1})$。于是，对任何 $r$，$\operatorname{Sing}_{\mathrm{sch}}(Y)$ 的长度都不小于 $r-1$，矛盾。因此 $Y$ 非奇异，且由 $\phi\in\C\{y_1,y_2,y_3,y_4^r\}$，曲面 $(\phi=y_4=0)$ 非奇异。于是 $E$ 非奇异，$k=3$ 的情形证完。

若 $k=2$，则不动点集是 $L=(y_3=y_4=0)$。于是 $E=L\ (\subset Y)$ 非奇异，且 $Y$ 的方程可以写成 $y_3f_3+y_4f_4=0$ 的形式。设 $\pi\colon B_LY\to Y$ 是沿 $L$ 爆破 $Y$ 所得的态射。$B_LY\subset B_L\C^4$，故 $\pi$ 的任意纤维的维数至多为 1。由于 $L$ 在原点以外是 Cartier 的，$\pi$ 在原点以外是同构。具体计算表明
\[
\pi^{-1}(0)\simeq\PP^1\iff\operatorname{mult}_0f_3\ge1\ \text{且}\ \operatorname{mult}_0f_4\ge1.
\]
因此，由引理 6.2，$X$ 在 0 处奇异 $\iff$ $E$ 不是 Q-Cartier 的。
\end{proof}

上面的证明中用到了下面这个容易证明的事实。

\begin{lemmap}{7.29}
设 $(R,M)$ 是带有有限交换群 $G$ 作用的局部环。此时，任意 $G$-不变理想 $I$ 都有由 $G$-特征向量组成的极小生成元集。
\end{lemmap}

\begin{proof}
设 $v\in I$ 使 $v+IM\in I/IM$ 是 $G$-特征向量，且 $\chi$ 为其

% [原文 p.272]
特征（character）。令 $w=\sum_{g\in G}\chi(g)^{-1}g(v)/|G|\in I$。此时，$h(w)=\chi(h)\cdot w$，且 $w\equiv v\bmod IM$ 成立。对 $I$ 的各个生成元应用这一操作即可。
\end{proof}

下面的推论完成 7.21 的阶段 5。

\begin{corollaryp}{7.30}
$B'^{(\ell)}_{tf}$ 与 $X'_{tf}$ 非奇异。
\end{corollaryp}

\begin{proof}
设 $E$ 是 $B'^{(\ell)}_{tf}$ 的不可约分支。此时，$T_E$ 固定 Q-Cartier 除子 $B'^{(\ell)}_{tf}$ 的每个点（7.21(3)）。$X'_{tf}$ 的奇点是至多指数 1 的终端奇点（命题 7.23(1)），从而孤立 cDV 奇点（推论 5.38）。因此，应用引理 7.28，便得到 $B'^{(\ell)}_{tf}$ 的每个点处 $B'^{(\ell)}_{tf}$ 与 $X'_{tf}$ 的非奇异性。$X'_{tf}$ 在 $\bigcup_\ell B'^{(\ell)}_{tf}$ 之外非奇异，故证明完成。
\end{proof}

那么，完成 7.21 的阶段 6。这样一来，定理 7.19 的证明也就完成了。

\begin{lemmap}{7.31}
$B'_{tf}$ 在局部解析意义下是 nc 除子。
\end{lemmap}

\begin{proof}
任取点 $P\in B'_{tf}$。设 $E_i$ 是 $B'_{tf}$ 的不可约分支（$\ni P$）。设 $g_i\in T_{E_i}$ 是其在 $X'_{tf}$ 上的作用不平凡的元素。设 $G$ 是由这些 $g_i$ 生成的子群。由引理 7.29，$E_i$ 是芽 $(P\in X'_{tf})$ 的坐标超平面。
\end{proof}

\section{特殊半稳定翻转}
\takeaway{本节证明三维特殊半稳定翻转的存在性：在 dlt、$S$ 的某分支与极射线负交等假设下，借由 7.34 造出 lc 补除子，把翻转归约到三维标准翻模的存在。}

本节目的是证明三维特殊半稳定翻转的存在性。这一结果及其证明与 [K+92, 20.8] 密切相关，而那里的证明又可追溯到 [Sho92]。

这一结果有几个假设略有不同的变体。条件 7.32(3) 是特殊翻转的条件，而条件 7.32(2) 是由下一节所需的主张决定的。

% [Tutor Note: 「特殊」体现在三条假设：S 是 Cartier 且 f-平凡、S 的某分支与极射线负交、$\lfloor B_X\rfloor$ 之外终端；这些正是下一节证明所需。]
\begin{theorem*}{7.32}
（特殊半稳定翻转的存在）设 $(X,S+B_X)$ 是 Q-分解三维对，$f\colon X\to Y$ 是射线型 $(K+B_X)$-翻转收缩，$S$ 是 $f$-平凡的有效 Cartier 除子，$B_X$ 是有效 Q-除子，记 $C=\operatorname{Ex}(f)$。假设：
% [原文 p.273]
\begin{enumerate}[label=(\arabic*),leftmargin=2.4em,itemsep=1pt,topsep=2pt,parsep=0pt]
\item $(X,S+B_X)$ 是 dlt 的。
\item $(X\setminus(S\cup\lfloor B_X\rfloor),0)$ 是终端的。
\item $S$ 含有满足 $(S_1\cdot C)<0$ 的不可约分支 $S_1$。
\end{enumerate}
此时，$f$ 有翻转。
\end{theorem*}

\begin{remark}{7.33}
上述定理在射影的、非固有代数的以及解析的设定下都同样成立。本节的论证在代数或解析设定下都能奏效。（若只想在射影簇的 MMP 范围内证明，参见注意 7.37。）
\end{remark}

证明的关键阶段是下面叙述的命题 7.34。首先说明由该命题可得定理 7.32，最后证明命题 7.34。

\begin{proposition*}{7.34}
设 $(X,S_1+S_2+B)$ 是 Q-分解三维对，$S_1,S_2$ 不可约，且 $f\colon X\to Y$ 是射线型 $(K+S_1+S_2+B)$-翻转收缩。假设：
\begin{enumerate}[label=(\arabic*),leftmargin=2.4em,itemsep=1pt,topsep=2pt,parsep=0pt]
\item $(X,S_1+S_2+B)$ 是 dlt 的。
\item $\lfloor B\rfloor=0$。
\item $(S_1\cdot C)<0$ 且 $(S_2\cdot C)>0$，其中 $C=\operatorname{Ex}(f)$。
\item $S_1$ 在 $X$ 上余维 2 处是 Cartier 的。
\end{enumerate}
此时，把 $X$ 收缩到 $C$ 的充分小邻域，存在 Weil 除子 $D^+\subset X$，使 $(X,S_1+S_2+D^++\lfloor 2B\rfloor)$ 是 lc 的，且 $K_X+S_1+S_2+D^++\lfloor 2B\rfloor\sim0$。
\end{proposition*}

\begin{remark}{7.35}
$D^+$ 与 [Sho92] 中引入的 1-补足（1-complement）密切相关。读本证明可知，自然对象并不是 $D^+$，而是位于适当奇点消解上的另一个除子（记为 $D^\sharp$）。
\end{remark}

\par\medskip\noindent\textbf{（命题 7.34）$\Rightarrow$（定理 7.32）的证明}\quad 由 $(S_1\cdot C)<0$，有 $C\subset S_1$。由于 $(S\cdot C)=0$，存在不可约分支 $S_2\subset S$ 使 $(S_2\cdot C)>0$。特别地，注意 $S_2$ 是 $f$-丰富的。对适当的 $0<\delta\ll1$，记
\[
S_3:=S-S_1-S_2,\qquad B:=(1-\delta)(S_3+B_X).
\]
由于 $S$ 是约化 Cartier 除子且 $(X,S+B_X)$ 是 dlt 的，$S_1$ 在 $X$ 上余维 2 处是 Cartier 的（推论 5.55）。命题 7.34 的条件得到满足，故把 $X$ 收缩到 $C$ 的周围之后，得到除子 $D^+\subset X$。

% [原文 p.274]
此后不再假设 $X$ 是 Q-分解的或 $f$ 是射线型的。因此，由推论 6.7，把 $X$ 收缩到 $C$ 的周围对翻转存在性的证明没有妨碍。今后也作同样的收缩。注意
\[
\operatorname{Supp}(S+\lfloor B_X\rfloor)\subset\operatorname{Supp}(S_1+S_2+\lfloor 2B\rfloor).
\]
由于 $(S_1\cdot C)<0$ 且 $(S_2\cdot C)>0$，存在 $f$-平凡的 Cartier 除子 $T$，它是 $\operatorname{Supp}(S_1+S_2+\lfloor 2B\rfloor)$ 的成分的正系数线性组合。收缩 $Y$，可假设 $T\sim0$。设 $m$ 是 $T$ 的系数的最小公倍数。设 $\pi\colon X^*\to X$ 是由取 $T$ 的 $m$ 次根所得的循环覆盖的不可约分支（定义 2.50）。此时，$T^*:=\pi^*(T)/m$ 是约化 Cartier 除子，且 $T^*\sim0$（命题 7.23(1)）。$X^*\setminus T^*$ 非奇异（定理 7.32(2)）。记 $D':=\pi^*(D^+)$，则
\[
K_{X^*}+T^*+D'=\pi^*(K_X+S_1+S_2+D^++\lfloor 2B\rfloor)
\]
成立（这是关于余维 1 的主张，故由命题 7.23(3) 可知）。于是，$(X^*,T^*+D')$ 是 lc 的（命题 5.20）。

% [Tutor Note: 此处「定理 7.32(2)」疑为命题 7.23(2) 之误：7.32(2) 讲终端性，而 $X^*\setminus T^*$ 的非奇异性正来自 7.23(2) 的 $X'_n\setminus D'_n$。]

\par\medskip\noindent\textbf{主张 7.36}\quad 设 $B^*\in|2D'|$ 是一般元。此时，$B^*$ 是约化的，$(X^*,(1/2)B^*)$ 是标准的，并且进一步 $2K_{X^*}+B^*\sim0$ 成立。

\begin{proof}
由于 $X^*\setminus T^*$ 非奇异，收缩 $Y$ 之后，完全线性系 $|2D'|$ 在 $T^*$ 之外没有基点。由此，由引理 5.17，$(X^*\setminus T^*,B^*)$ 是标准的。由推论 2.33，$(X^*,T^*+(1/2)B^*)$ 是 lc 的。

为确认 $(X^*,(1/2)B^*)$ 是标准的，设 $E$ 是某个双有理态射 $h\colon Z\to X^*$ 中出现的例外除子。若 $h(E)\not\subset\operatorname{Supp}T^*$，则由 $(X^*\setminus T^*,B^*)$ 是标准的，$a(E,X^*,(1/2)B^*)\ge0$。反之，若 $h(E)\subset\operatorname{Supp}T^*$，则
\[
a(E,X^*,(1/2)B^*)\ge1+a(E,X^*,T^*+(1/2)B^*)\ge1-1=0.
\]
\end{proof}

$(X^*,(1/2)B^*)$ 满足引理 6.11 的假设。因此，$f\colon X\to Y$ 的翻转的存在性归约到三维标准翻模（flop）的存在性（推论 6.9）。由推论 6.45，三维标准翻模存在，故定理 7.32 得证。\hfill$\blacksquare$

\begin{remark}{7.37}
在定理 7.32 中，若 $Y$ 是射影三维簇，则本节的论证只需稍作修改（不必移到开部分集合），就能成为关于射影簇的论证。

% [原文 p.275]
在命题 7.34 中，加上 $Y$ 是射影的且具有丰富除子 $L$ 的假设，并把结论改为：
对 $\nu\gg0$，存在 Weil 除子 $D^+\subset X$，使 $(X,S_1+S_2+D^+)$ 是 lc 的，且 $K_X+S_1+S_2+D^+\sim\nu f^*L$。

于是，代替 $K_X+S_1+S_2+D^+\sim0$，有 $K_X+S_1+S_2+D^+\sim_f0$ 成立。在上述证明中，令 $\nu\equiv0\bmod(m)$，把一般元 $V\sim\nu f^*L-T$ 选得既不与 $C$ 相交、又与 $T$ 横截地相交。然后取 $T+V\sim\nu f^*L$ 的 $m$ 次根。于是，由定义 2.49，可以构成全局分歧覆盖，且（命题 7.34）$\Rightarrow$（定理 7.32）的论证能够奏效。关于修改后的命题 7.34，只要把命题 7.38 同样修改，证明便可原样通用。
\end{remark}

% [Tutor Note: 注意 7.37 说明射影情形只需把结论放宽到 $K_X+S_1+S_2+D^+\sim\nu f^*L$，用全局分歧覆盖代替局部循环覆盖。]
\par\medskip\noindent\textbf{命题 7.34 的证明}\quad 设 $g\colon X'\to X$ 是 $(X,S_1+S_2+B)$ 的对数奇点消解，记 $S:=g_*^{-1}S_1$。注意由命题 5.61，$(S_1,(S_2+B)|_{S_1})$ 是 plt 的。因此，$D:=(g|_S)_*^{-1}(S_2|_{S_1})$ 是 $S$ 上的约化除子，且由
\[
K_S+D+\Delta\equiv(g|_S)^*(K_{S_1}+(S_2+B)|_{S_1})
\]
定义的 $S$ 上的 Q-除子 $\Delta$ 满足 $\lfloor\Delta\rfloor\le0$。又 $(g|_S)_*\Delta=B|_{S_1}\ge0$ 是显然的，故由引理 2.30，$(S,D+\Delta)$ 是 plt 的。

即使把 $B$ 的系数稍作增加，不仅关于 $\lfloor\Delta\rfloor$ 的主张不变，命题 7.34 的假设与结论也都不变（推论 2.39）。因此，通过这样做，可以使 $B,\Delta$ 的系数都不属于 $(1/2)\Z$。（这个 $(1/2)\Z$ 正是取整函数 $x\mapsto\lfloor 2x\rfloor$ 的不连续点全体。）
\[
\begin{array}{ccc}
X' & \supset & S\ \supset\ D\\[2pt]
{\scriptstyle f\circ g}\big\downarrow & & {\scriptstyle p}\big\downarrow\quad\big\downarrow\\[2pt]
Y & \supset & f(S_1)\xrightarrow{\ \text{正规化}\ }T\ \supset\ p(D)
\end{array}
\]

首先，把命题 7.38 应用于 $T:=f(S_1)$ 的正规化、$p:=f\circ g|_S$，得到使 $(S,D+\lfloor 2\Delta\rfloor+D^\sharp_S)$ 为 lc 的除子 $D^\sharp_S\in|-K_S-D-\lfloor 2\Delta\rfloor|$。下面把 $D^\sharp_S$ 延拓到 $X'$ 上。取 $0<\eps\ll1$ 且 $\eps\in\Q$。此时，$-(K_X+S_1+(1-\eps)(S_2+B))$ 是 $f$-丰富的。设 $B'_\eps$ 为由
\[
K_{X'}+S+B'_\eps\equiv g^*(K_X+S_1+(1-\eps)(S_2+B))
\]
确定的 $X'$ 上的 Q-除子。

% [原文 p.276]
当 $\eps\searrow0$ 时，$B'_\eps$ 是除子的单调增加列，且进一步注意 $\lim g_*B'_\eps=S_2+B$ 与 $\lim B'_\eps|_S=D+\Delta$ 成立。由 $(X,S_1+(1-\eps)(S_2+B))$ 是 plt 的（命题 5.51），$\lfloor B'_\eps\rfloor\le0$。由此可知 $\lfloor B'_\eps|_S\rfloor\le0$ 且 $\lfloor 2B'_\eps|_S\rfloor=D+\lfloor 2\Delta\rfloor$。由命题 7.39，$D^\sharp_S\in|-K_S-D-\lfloor 2\Delta\rfloor|$ 被延拓为 $D^\sharp\in|-K_{X'}-S-\lfloor 2B'_\eps\rfloor|$。

最后，确认 $D^+:=g_*D^\sharp$ 满足命题 7.34 的条件。由 $\lfloor B'_\eps\rfloor\le0$ 与 $\lim g_*B'_\eps=S_2+B$，可知 $\lfloor 2g_*B'_\eps\rfloor=S_2+\lfloor 2B\rfloor$。因此，$K_X+S_1+S_2+D^++\lfloor 2B\rfloor\sim0$ 由 $K_{X'}+S+\lfloor 2B'_\eps\rfloor+D^\sharp\sim0$ 导出。进一步，
\[
\begin{aligned}
&(X,S_1+S_2+D^++\lfloor 2B\rfloor)\ \text{是 lc}\iff(X',S+\lfloor 2B'_\eps\rfloor+D^\sharp)\ \text{是 lc}\\
&\iff(S,(\lfloor 2B'_\eps\rfloor+D^\sharp)|_S)=(S,D+\lfloor 2\Delta\rfloor+D^\sharp_S)\ \text{是 lc}.
\end{aligned}
\]
其中，第一个等价性由引理 2.30，第二个由定理 5.50(2)。最后一对是 lc 的，这已在第 1 阶段确认。\hfill$\blacksquare$

% [Tutor Note: 命题 7.38 与 7.39 是曲面上的补除子工具：先在 $S$ 上造出使对 lc 的除子，再借对数消解把它提升回三维的 $D^+$。]
\begin{propositionp}{7.38}
设 $T$ 是正规曲面，$p\colon S\to T$ 是奇点消解，$C$ 为例外集。设 $D\subset S$ 是 Weil 除子，$\Delta$ 是（不必有效的）Q-除子，满足以下条件：
\begin{enumerate}[label=(\arabic*),leftmargin=2.4em,itemsep=1pt,topsep=2pt,parsep=0pt]
\item $\operatorname{Supp}D$ 与 $\operatorname{Supp}\Delta$ 没有公共不可约分支。
\item $(S,D+\Delta)$ 是 plt 的，且 $D+\Delta$ 是 snc 除子。
\item 在 $p(C)$ 的各点的邻域中 $p_*D\ne0$。
\item $\lfloor\Delta\rfloor\le0$，且 $p_*\Delta$ 是有效的。
\item $-(K_S+D+\Delta)$ 是 $p$-nef 的。
\end{enumerate}
此时，把 $S$ 收缩到 $C$ 的周围，存在 Weil 除子 $D^\sharp_S\subset S$，使 $(S,D+\lfloor 2\Delta\rfloor+D^\sharp_S)$ 是 lc 的，且 $K_S+D+\lfloor 2\Delta\rfloor+D^\sharp_S\sim0$。
\end{propositionp}

\begin{proof}
由命题 5.51，$D$ 非奇异。任取点 $t\in p(C)$ 与充分小的邻域 $U\ni t$。由引理 3.38，$(T,p_*(D+\Delta))$ 是 plt 的，故 $p(D)|_U$ 非奇异（命题 5.51），而且非空（命题 7.38(3)）。也就是说，$D\cap p^{-1}(U)$ 恰好有一个不含于 $C$ 的不可约分支。另一方面，$D\cap p^{-1}(U)$ 是连通的（推论 5.49）。因此 $D\cap p^{-1}(U)$ 在 $p^{-1}(t)$ 中恰有一个点（记为 $x_t$）。由于 $t\in p(C)$ 是任意的，$p\colon D\to T$ 是有限的。
\end{proof}